\documentclass[11pt,a4paper]{article}

% Preambolo condiviso dai documenti del progetto KeyItaly.
% Stile accademico sobrio: nessun colore decorativo, tipografia standard.

\usepackage[T1]{fontenc}
\usepackage[utf8]{inputenc}
\usepackage[italian]{babel}
\usepackage{lmodern}
\usepackage{microtype}

\usepackage[a4paper,margin=2.8cm]{geometry}
\usepackage{graphicx}
\usepackage{booktabs}
\usepackage{array}
\usepackage{enumitem}
\usepackage{xcolor}
\usepackage{tikz}
\usepackage{float}
\usepackage{pdflscape}
\usepackage[font=small,labelfont=bf,width=.92\textwidth]{caption}
\usepackage{listings}

\usepackage[hidelinks,pdfusetitle]{hyperref}
\usepackage{url}
\urlstyle{same}

\usetikzlibrary{positioning,shapes.geometric,arrows.meta,calc,fit,backgrounds}

% Interlinea leggermente aperta, più leggibile in stampa
\usepackage{setspace}
\onehalfspacing

% Elenchi più compatti dei default di LaTeX
\setlist[itemize]{itemsep=2pt,topsep=4pt,parsep=0pt}
\setlist[enumerate]{itemsep=2pt,topsep=4pt,parsep=0pt}
\setlist[description]{itemsep=4pt,topsep=6pt}

% ---------------------------------------------------------------------------
% Diagrammi
%
% La codifica cromatica è la stessa in tutto il documento:
%   grigio = accessibile a chiunque
%   azzurro = riservato all'utente registrato
%   rosso = riservato all'amministratore
%   giallo = vista JSP
%   verde = pattern di URL protetto da un filtro
% ---------------------------------------------------------------------------
\definecolor{cPubblico}{HTML}{ECECEC}
\definecolor{cUtente}{HTML}{CFE2F7}
\definecolor{cAdmin}{HTML}{F7D2D2}
\definecolor{cVista}{HTML}{FDEDC8}
\definecolor{cFiltro}{HTML}{D8EBD5}
\definecolor{cBordo}{HTML}{5A5A5A}

\tikzset{
  nodo/.style={
    draw=cBordo, line width=.5pt, rounded corners=2pt,
    align=center, font=\scriptsize, inner sep=2pt,
    minimum height=7mm, text width=20mm},
  pubblico/.style={nodo, fill=cPubblico},
  utente/.style={nodo, fill=cUtente},
  admin/.style={nodo, fill=cAdmin},
  vista/.style={nodo, fill=cVista},
  filtro/.style={nodo, fill=cFiltro},
  linea/.style={draw=cBordo, line width=.5pt},
  freccia/.style={linea, -{Stealth[length=4pt,width=3pt]}},
  tratteggio/.style={freccia, dashed, dash pattern=on 2.5pt off 2pt},
  % Notazione entità-relazione
  entita/.style={draw=black, line width=.6pt, fill=white,
    align=center, font=\small, minimum height=8mm, inner xsep=8pt},
  relazione/.style={draw=black, line width=.6pt, fill=white, diamond,
    aspect=2.2, align=center, font=\scriptsize, inner sep=1pt},
  composto/.style={draw=black, line width=.6pt, fill=white, ellipse,
    align=center, font=\scriptsize, inner sep=2pt},
  pallino/.style={draw=black, line width=.5pt, fill=white, circle,
    inner sep=0pt, minimum size=3.2pt},
  pallinoPieno/.style={pallino, fill=black},
  etichetta/.style={font=\scriptsize, inner sep=1.5pt},
}

% Attributo di un'entità nella notazione entità-relazione.
%   #1 stile del pallino (pallino oppure pallinoPieno)
%   #2 punto di attacco sul bordo dell'entità
%   #3 spostamento fino al pallino
%   #4 ancoraggio dell'etichetta rispetto al pallino
%   #5 testo dell'etichetta
\newcommand{\attributo}[5]{%
  \draw[line width=.5pt] (#2) -- ++(#3) node[#1] (ultimoAttr) {};
  \node[etichetta, anchor=#4, align=center] at (ultimoAttr) {#5};%
}

% Le figure occupano quasi tutta la larghezza del testo
\newcommand{\figuraLarga}[3]{%
  \begin{figure}[H]
    \centering
    \includegraphics[width=\linewidth]{#1}
    \caption{#2}\label{#3}
  \end{figure}}

% Stile per i frammenti di codice e configurazione
\lstdefinestyle{sobrio}{
  basicstyle=\ttfamily\footnotesize,
  keywordstyle=\bfseries,
  commentstyle=\itshape,
  numbers=none,
  frame=single,
  rulecolor=\color{black!35},
  framesep=6pt,
  breaklines=true,
  showstringspaces=false,
  captionpos=b,
  extendedchars=true,
  literate={à}{{\`a}}1 {è}{{\`e}}1 {é}{{\'e}}1 {ì}{{\`i}}1 {ò}{{\`o}}1 {ù}{{\`u}}1
}
\lstset{style=sobrio}

% Frontespizio
\newcommand{\frontespizio}[2]{%
  \begin{titlepage}
    \centering
    \vspace*{2cm}
    {\LARGE Università degli Studi di Salerno\par}
    \vspace{4mm}
    {\large Dipartimento di Informatica\par}
    {\large Corso di Laurea in Informatica\par}
    \vspace{2.2cm}
    {\large Tecnologie Software per il Web\par}
    \vspace{1.4cm}
    \rule{\linewidth}{0.4pt}\par
    \vspace{6mm}
    {\huge\bfseries #1\par}
    \vspace{4mm}
    {\Large\itshape KeyItaly\par}
    \vspace{3mm}
    {\large #2\par}
    \vspace{6mm}
    \rule{\linewidth}{0.4pt}\par
    \vfill
    \begin{minipage}[t]{0.45\textwidth}
      \raggedright
      \textbf{Docente}\par
      \vspace{2mm}
      Prof.ssa Rita Francese
    \end{minipage}%
    \begin{minipage}[t]{0.45\textwidth}
      \raggedleft
      \textbf{Studenti}\par
      \vspace{2mm}
      Giovanni Cerchia\quad 0512120818\par
      Andrea Compagnone\quad 0512116159
    \end{minipage}
    \vspace{1.4cm}
    \par
    {\large Anno Accademico 2023/2024\par}
    \vspace{1cm}
  \end{titlepage}}

% ===========================================================================
% Diagrammi del progetto KeyItaly, disegnati con TikZ.
% Gli stili condivisi sono definiti in preambolo.tex.
% ===========================================================================

% Entità dello schema concettuale: intestazione e, sotto la riga, gli attributi.
\newcommand{\boxEntita}[2]{%
  \begin{tabular}{@{}c@{}}
    \textbf{#1}\\[1.5pt]
    \hline
    \rule{0pt}{2.6ex}%
    \begin{tabular}{@{}l@{}}#2\end{tabular}
  \end{tabular}}

% ---------------------------------------------------------------------------
% Diagramma navigazionale
% ---------------------------------------------------------------------------
\newcommand{\diagrammaNavigazionale}{%
\begin{tikzpicture}[x=1cm, y=1cm, every node/.append style={font=\small}]

  \node[pubblico] (home) at (8,0) {Home};

  \node[pubblico] (carrello)  at ( 0,-2) {Carrello};
  \node[pubblico] (categorie) at ( 4,-2) {Categorie\\catalogo};
  \node[pubblico] (catalogo)  at ( 8,-2) {Catalogo};
  \node[pubblico] (accesso)   at (12,-2) {Login /\\Registrazione};
  \node[admin]    (listaord)  at (16,-2) {Lista ordini};

  \draw[linea] (home.south) -- (8,-1.15);
  \draw[linea] (0,-1.15) -- (16,-1.15);
  \foreach \x in {0,4,8,12,16} { \draw[linea] (\x,-1.15) -- (\x,-1.65); }

  % Carrello
  \node[utente] (checkout)  at (0,-3.8) {Checkout};
  \node[utente] (riepilogo) at (0,-5.5) {Riepilogo\\ordine};
  \draw[linea] (carrello.south) -- (checkout.north);
  \draw[linea] (checkout.south) -- (riepilogo.north);

  % Categorie del catalogo
  \node[pubblico] at (5.3,-3.8) {Tastiere};
  \node[pubblico] at (5.3,-5.3) {Switch};
  \node[pubblico] at (5.3,-6.8) {Copritasti};
  \draw[linea] (categorie.south) -- (4,-6.8);
  \foreach \y in {-3.8,-5.3,-6.8} { \draw[linea] (4,\y) -- (4.25,\y); }

  % Catalogo
  \node[pubblico] at (9.3,-3.8) {Descrizione\\prodotto};
  \node[admin]    at (9.3,-5.3) {Cancella\\prodotto};
  \node[admin]    at (9.3,-6.8) {Aggiungi\\prodotto};
  \node[admin]    at (9.3,-8.3) {Modifica\\prodotto};
  \draw[linea] (catalogo.south) -- (8,-8.3);
  \foreach \y in {-3.8,-5.3,-6.8,-8.3} { \draw[linea] (8,\y) -- (8.25,\y); }

  % Accesso e profilo
  \node[utente] (profilo) at (12,-3.8) {Profilo};
  \draw[linea] (accesso.south) -- (profilo.north);
  \node[utente] at (13.3,-5.5)  {Modifica\\indirizzo};
  \node[utente] at (13.3,-7.1)  {Modifica metodo\\pagamento};
  \node[utente] at (13.3,-8.7)  {Mostra dettagli\\ordine};
  \node[utente] at (13.3,-10.3) {Stampa fattura\\ordine};
  \draw[linea] (profilo.south) -- (12,-10.3);
  \foreach \y in {-5.5,-7.1,-8.7,-10.3} { \draw[linea] (12,\y) -- (12.25,\y); }

  % Legenda
  \begin{scope}[shift={(0,-12.1)}]
    \node[pubblico, minimum height=5mm, text width=17mm] at (0,0)   {pubblico};
    \node[utente,   minimum height=5mm, text width=25mm] at (3.1,0) {utente registrato};
    \node[admin,    minimum height=5mm, text width=25mm] at (6.6,0) {amministratore};
  \end{scope}
\end{tikzpicture}}

% ---------------------------------------------------------------------------
% Mappa dei contenuti
% ---------------------------------------------------------------------------
\newcommand{\mappaContenuti}{%
\begin{tikzpicture}[x=1cm, y=1cm, every node/.append style={font=\small}]

  \node[pubblico] (home) at (7.5,0) {Home};

  \node[pubblico] (carrello)   at ( 0,-2.4)   {Carrello};
  \node[pubblico] (tastiere)   at ( 2.7,-2.4) {Tastiere};
  \node[pubblico] (switchp)    at ( 5.4,-2.4) {Switch};
  \node[pubblico] (copritasti) at ( 8.1,-2.4) {Copritasti};
  \node[utente]   (profilo)    at (11.0,-2.4) {Profilo};
  \node[admin]    (listaord)   at (14.5,-2.4) {Lista ordini};

  \foreach \n in {carrello,tastiere,switchp,copritasti,profilo,listaord}
    { \draw[freccia] (home) to[out=-90,in=90] (\n); }

  \node[utente]   (checkout)  at ( 0,-4.8)   {Checkout};
  \node[pubblico] (prodotto)  at ( 5.4,-4.8) {Prodotto};
  \node[utente]   (indirizzo) at ( 9.6,-4.8) {Indirizzo};
  \node[utente]   (pagamento) at (12.7,-4.8) {Metodo di\\pagamento};
  \node[utente]   (storico)   at (11.0,-7.1) {Storico\\ordini};
  \node[utente]   (ordine)    at ( 9.6,-9.4) {Ordine};
  \node[utente]   (fattura)   at (12.7,-9.4) {Fattura};

  \draw[freccia] (carrello) -- (checkout);
  \foreach \n in {tastiere,switchp,copritasti}
    { \draw[freccia] (\n) to[out=-90,in=90] (prodotto); }
  \draw[freccia] (profilo) to[out=-90,in=90] (indirizzo);
  \draw[freccia] (profilo) to[out=-90,in=90] (pagamento);
  \draw[freccia] (profilo) -- (storico);
  \draw[freccia] (storico) to[out=-90,in=90] (ordine);
  \draw[freccia] (ordine) -- (fattura);
  % L'amministratore raggiunge lo stesso dettaglio d'ordine dalla lista completa
  \draw[freccia] (listaord.south) -- (14.5,-10.7) -- (9.6,-10.7) -- (ordine.south);
\end{tikzpicture}}

% ---------------------------------------------------------------------------
% Schema concettuale della base di dati
% ---------------------------------------------------------------------------
\newcommand{\schemaER}{%
\begin{tikzpicture}[x=1cm, y=1cm,
    ent/.style={draw=black, line width=.6pt, fill=white, rounded corners=1pt,
                inner sep=5pt, font=\scriptsize, align=center},
    rel/.style={draw=black, line width=.6pt, fill=white, diamond, aspect=3.2,
                inner sep=1pt, minimum width=2.9cm, font=\scriptsize, align=center},
    card/.style={font=\scriptsize, inner sep=2pt, fill=white}]

  % --- Utente e specializzazioni ---
  \node[ent] (utente) at (6,0) {\boxEntita{Utente}{%
      \underline{username}\\ password\\ nome, cognome\\ email \textit{(unico)}\\ data di nascita}};

  \node[ent] (cliente) at (2.6,-4.6) {\boxEntita{Cliente}{%
      \textit{metodo di pagamento:}\\
      \quad nome, cognome carta\\
      \quad numero, scadenza, CVV\\[2pt]
      \textit{indirizzo di fatturazione:}\\
      \quad via, CAP, città}};

  \node[ent] (admin) at (9.6,-4.6) {\boxEntita{Admin}{\textit{nessun attributo}\\ \textit{proprio}}};

  \draw[line width=.6pt] (cliente.north) -- (2.6,-3.0) -- (9.6,-3.0) -- (admin.north);
  \draw[line width=.6pt] (6,-3.0) -- (6,-2.3);
  \draw[-{Triangle[length=5pt,width=7pt]}, line width=.9pt] (6,-2.6) -- (utente.south);
  \node[font=\scriptsize, anchor=west] at (6.25,-2.75) {generalizzazione};

  % --- Ordini ---
  \node[rel] (esecuzione) at (2.6,-7.6) {Esecuzione};
  \draw[line width=.6pt] (cliente.south) -- (esecuzione.north);
  \node[card, anchor=west] at (2.7,-6.9) {(0,N)};

  \node[ent] (ordine) at (2.6,-10.8) {\boxEntita{Ordine}{%
      \underline{ID}\\ prezzo totale\\ data ordine, data consegna\\
      destinatario: nome, cognome\\ consegna: via, CAP, città}};
  \draw[line width=.6pt] (esecuzione.south) -- (ordine.north);
  \node[card, anchor=west] at (2.7,-8.5) {(1,1)};

  % --- Acquisti ---
  \node[rel] (acquisto) at (7.6,-10.8) {Acquisto};
  \draw[line width=.6pt] (ordine.east) -- (acquisto.west);
  \node[card, anchor=south] at (5.5,-10.7) {(1,N)};
  \node[ent, draw=black!55, dash pattern=on 2pt off 1.5pt] (attrAcq) at (7.6,-8.6)
    {\begin{tabular}{@{}l@{}}quantità\\ prezzo di acquisto\\ IVA di acquisto\end{tabular}};
  \draw[line width=.5pt, draw=black!55, dash pattern=on 2pt off 1.5pt]
    (acquisto.north) -- (attrAcq.south);

  \node[ent] (prodotto) at (12.3,-10.8) {\boxEntita{Prodotto}{%
      \underline{ID}\\ nome, descrizione\\ prezzo, IVA\\ grafica\\ quantità in stock}};
  \draw[line width=.6pt] (acquisto.east) -- (prodotto.west);
  \node[card, anchor=south] at (10.0,-10.7) {(0,N)};

  % --- Specializzazioni del prodotto ---
  \node[ent] (tastiere)   at ( 9.0,-14.4) {Tastiere};
  \node[ent] (switchp)    at (12.3,-14.4) {Switch};
  \node[ent] (copritasti) at (15.4,-14.4) {Copritasti};
  \draw[line width=.6pt] (tastiere.north) -- (9.0,-13.5) -- (15.4,-13.5) -- (copritasti.north);
  \draw[line width=.6pt] (switchp.north) -- (12.3,-13.5);
  \draw[-{Triangle[length=5pt,width=7pt]}, line width=.9pt] (12.3,-13.5) -- (prodotto.south);
\end{tikzpicture}}

% ---------------------------------------------------------------------------
% Wireframe della pagina principale
% ---------------------------------------------------------------------------
\newcommand{\wireframeHome}{%
\begin{tikzpicture}[x=1cm, y=1cm,
    riquadro/.style={draw=black, line width=.7pt, align=center}]

  \def\W{12}

  \draw[riquadro] (0,0) rectangle (\W,-0.9);
  \foreach \x in {2,4,7.2,9.4} { \draw[line width=.7pt] (\x,0) -- (\x,-0.9); }
  \node[font=\small] at (1,-0.45) {Home};
  \node[font=\small] at (3,-0.45) {Catalogo};
  \node[font=\small, align=center] at (5.6,-0.45) {Barra\\di ricerca};
  \node[font=\small] at (8.3,-0.45) {Utente};
  \node[font=\small] at (10.7,-0.45) {Carrello};

  \draw[riquadro] (0,-0.9) rectangle (\W,-3.6);
  \node[font=\LARGE] at (6,-2.25) {News};
  \node[font=\large] at (0.7,-2.25) {$\Leftarrow$};
  \node[font=\large] at (11.3,-2.25) {$\Rightarrow$};

  \draw[riquadro] (0,-3.6) rectangle (\W,-5.6);
  \foreach \x in {4,8} { \draw[line width=.7pt] (\x,-3.6) -- (\x,-5.6); }
  \node[font=\large] at (2,-4.6)  {Tastiere};
  \node[font=\large] at (6,-4.6)  {Copritasti};
  \node[font=\large] at (10,-4.6) {Switch};

  \draw[riquadro] (0,-5.6) rectangle (\W,-8.4);
  \node[font=\small] at (6,-6.0) {I nostri prodotti};
  \foreach \x in {0.9,3.6,6.3,9.0}
    { \draw[line width=.7pt] (\x,-6.4) rectangle ++(2.1,-1.4); }
  \node[font=\small] at (5.2,-8.1) {$\leftarrow$};
  \node[font=\small] at (6.8,-8.1) {$\rightarrow$};

  \draw[riquadro] (0,-8.4) rectangle (\W,-9.6);
  \node[font=\large] at (6,-9.0) {Footer};
\end{tikzpicture}}

% ---------------------------------------------------------------------------
% Palette cromatica
% ---------------------------------------------------------------------------
\newcommand{\paletteColori}{%
\begin{tikzpicture}[x=1mm,y=1mm]
  \definecolor{deepskyblue}{HTML}{007BFF}
  \definecolor{blueeyes}{HTML}{1663C7}
  \definecolor{snowdrift}{HTML}{F8F9FA}
  \definecolor{dune}{HTML}{333333}

  \foreach \i/\col/\hex/\nome in {
    0/deepskyblue/007BFF/Deep Sky Blue,
    1/blueeyes/1663C7/Blue Eyes,
    2/snowdrift/F8F9FA/Snow Drift,
    3/white/FFFFFF/White,
    4/dune/333333/Dune}
  {
    \fill[\col,draw=black!45] (\i*28,0) rectangle ++(25,25);
    \node[anchor=north,font=\footnotesize\ttfamily] at (\i*28+12.5,-2) {\#\hex};
    \node[anchor=north,font=\scriptsize,text width=26mm,align=center]
      at (\i*28+12.5,-7) {\nome};
  }
\end{tikzpicture}}

% ---------------------------------------------------------------------------

% ---------------------------------------------------------------------------
% Mappa delle servlet, dei filtri e delle viste
% ---------------------------------------------------------------------------
\newcommand{\mappaServlet}{%
\begin{tikzpicture}[x=1cm, y=1cm,
    srv/.style={nodo, fill=cUtente, text width=40mm, font=\footnotesize\ttfamily},
    jsp/.style={nodo, fill=cVista,  text width=40mm, font=\footnotesize\ttfamily},
    esito/.style={nodo, fill=cAdmin, text width=40mm, font=\footnotesize},
    tit/.style={font=\small\bfseries, anchor=west},
    flt/.style={font=\footnotesize, anchor=west}]

  % ======================= Area pubblica =======================
  \node[tit] at (-2.2,1.2) {Area pubblica};

  \node[srv] (r1) at (0,0)      {/};
  \node[jsp] (v1) at (5.4,0)    {home.jsp};

  \node[srv] (r2) at (0,-1.9)   {/Catalogo\\/Tipo};
  \node[jsp] (v2) at (5.4,-1.9) {catalogo.jsp\\catalogoAdmin.jsp};

  \node[srv] (r3) at (0,-3.9)   {/DescrizioneProdotto};
  \node[jsp] (v3) at (5.4,-3.9) {descrizione.jsp};

  \node[srv] (r4) at (0,-6.1)   {/Login\\/Logout\\/Registrazione};
  \node[jsp] (v4) at (5.4,-6.1) {login.jsp};

  \node[srv] (r5) at (0,-8.3)   {/AggiungiProdotto\\/RimuoviProdotto};
  \node[jsp] (v5) at (5.4,-8.3) {carrello.jsp};

  \node[srv] (r6) at (0,-10.6)    {/ajax/search\\/CheckEmail\\/api/GetImage};
  \node[esito] (v6) at (5.4,-10.6) {risposte JSON\\e immagini del catalogo};

  \foreach \a/\b in {r1/v1, r2/v2, r3/v3, r4/v4, r5/v5, r6/v6}
    { \draw[freccia] (\a) -- (\b); }

  % ======================= Area /user/* =======================
  \node[tit] at (11.3,1.2) {Riservata a \texttt{/user/*}};

  \node[srv] (u1) at (13.5,-0.2) {/user/Profilo\\/user/ModificaIndirizzo\\/user/ModificaPagamento};
  \node[jsp] (w1) at (18.9,-0.2) {profilo.jsp};

  \node[srv] (u3) at (13.5,-2.6) {/user/Checkout};
  \node[jsp] (w3) at (18.9,-2.6) {checkout.jsp};
  \node[jsp] (w4) at (18.9,-4.0) {acquisto.jsp};

  \node[srv] (u4) at (13.5,-5.5)   {/user/StampaFattura};
  \node[esito] (w5) at (18.9,-5.5) {fattura in PDF};

  \draw[freccia] (u1) -- (w1);
  \draw[freccia] (u3) -- (w3);
  \draw[freccia] (w3) -- (w4);
  \draw[freccia] (u4) -- (w5);

  \node[draw=cBordo, dash pattern=on 3pt off 2pt, rounded corners=3pt,
        fit=(u1)(w1)(u4)(w5), inner sep=7pt] (boxuser) {};
  \node[flt] at (11.3,-7.0) {protetta da \texttt{UserFilter}};

  % ======================= Area /admin/* =======================
  \node[tit] at (11.3,-8.4) {Riservata a \texttt{/admin/*}};

  \node[srv] (a1) at (13.5,-9.6) {/admin/OrdiniAdmin};
  \node[jsp] (x1) at (18.9,-9.6) {ordiniAdmin.jsp};

  \node[srv] (a2) at (13.5,-11.1) {/admin/ModificaProdotto};
  \node[jsp] (x2) at (18.9,-11.1) {modifica.jsp};

  \node[srv] (a3) at (13.5,-12.9) {/admin/SaveProdotto\\/admin/DeleteProdotto};
  \node[jsp] (x3) at (18.9,-12.9) {catalogoAdmin.jsp};

  \foreach \a/\b in {a1/x1, a2/x2, a3/x3} { \draw[freccia] (\a) -- (\b); }

  \node[draw=cBordo, dash pattern=on 3pt off 2pt, rounded corners=3pt,
        fit=(a1)(x1)(a3)(x3), inner sep=7pt] (boxadmin) {};
  \node[flt] at (11.3,-14.4) {protetta da \texttt{AdminFilter}};

  % ======================= Filtro sulle viste =======================
  \node[srv, fill=cFiltro] (jspf) at (0,-12.5) {JspFilter};
  \node[esito] (forbidden) at (5.4,-12.5) {403 Forbidden};
  \draw[freccia] (jspf) -- (forbidden);
  \node[font=\footnotesize, text width=80mm, align=center] at (2.7,-13.9)
    {ogni accesso diretto a \texttt{/pages/views/*}};

\end{tikzpicture}}


\title{KeyItaly --- Project Proposal}
\author{Giovanni Cerchia \and Andrea Compagnone}
\date{Anno Accademico 2023/2024}

\begin{document}

% Il frontespizio e l'indice non ricevono ancoraggi di pagina: senza questo,
% la numerazione ripartendo da 1 creerebbe due destinazioni con lo stesso nome.
\hypersetup{pageanchor=false}
\frontespizio{Project Proposal}{Proposta di progetto e sistema realizzato}

\pagenumbering{roman}
\tableofcontents
\clearpage
\pagenumbering{arabic}
\hypersetup{pageanchor=true}

\section{Obiettivo del progetto}

\emph{KeyItaly} si propone come sito di commercio elettronico per la vendita di
tastiere meccaniche e di accessori quali switch e keycaps, rivolto sia agli
appassionati sia ai professionisti che cercano prodotti di qualità e
personalizzabili.

Il catalogo offre tastiere adatte a esigenze diverse --- dal gioco allo sviluppo
software, dal lavoro d'ufficio alla scrittura --- affiancate dai componenti
necessari a personalizzare una tastiera già assemblata. L'interfaccia è
progettata in modo responsive, così da restare utilizzabile da qualsiasi
dispositivo e a qualsiasi risoluzione.

\section{Analisi dei competitor}

L'analisi ha riguardato due negozi di riferimento del settore, uno
internazionale e uno italiano, dai quali sono state tratte le scelte di
struttura e di navigazione adottate nel progetto.

\subsection{Keychron}

\noindent\url{https://www.keychron.com/}

\figuraLarga{figures/keychron-home.jpeg}
  {Pagina principale di Keychron.}{fig:keychron}

Keychron è specializzato in tastiere meccaniche wireless. L'interfaccia è
essenziale: un banner centrale presenta novità e prossime uscite, mentre
l'header raccoglie i collegamenti al catalogo completo (\emph{All products}), a
quello delle sole tastiere (\emph{All keyboards}), alla barra di ricerca,
all'area utente e al carrello.

Il catalogo può essere filtrato per tipo, serie, layout, disponibilità e prezzo,
e ordinato per nome, data, prezzo, prodotti in evidenza e più venduti. Ogni
prodotto dispone di una scheda con descrizione, fotografie e selezione di
colore e tipo di switch. Il carrello è utilizzabile senza autenticazione, che
viene richiesta solo al momento del pagamento: è la soluzione adottata anche in
KeyItaly.

\subsection{CoffeeKeys}

\noindent\url{https://coffeekeys.eu/}

\figuraLarga{figures/coffeekeys-home.jpeg}
  {Pagina principale di CoffeeKeys.}{fig:coffeekeys-home}

CoffeeKeys è il primo negozio italiano dedicato esclusivamente alle tastiere
meccaniche. La struttura è simile a quella di Keychron: banner delle novità in
apertura, logo centrale affiancato da ricerca, area utente e carrello, e un menù
che articola il catalogo per brand e per dimensione, con le categorie
\emph{DIY} e \emph{Group Buys}. Completano il sito le sezioni
\emph{Accessori}, \emph{Contattaci} e \emph{FAQ}, e una sezione \emph{News} con
guide all'acquisto.

\figuraLarga{figures/coffeekeys-prodotto.jpeg}
  {Scheda di prodotto su CoffeeKeys.}{fig:coffeekeys-prodotto}

La scheda di prodotto raccoglie fotografie da più angolazioni, selezione di
colore e switch, specifiche tecniche e articoli correlati. È il modello a cui si
ispira la pagina di dettaglio di KeyItaly.

\section{Funzionalità del sito}

Le funzionalità sono organizzate su tre livelli di accesso.

\subsection{Utente non registrato}

\begin{itemize}
  \item consultare il catalogo completo o filtrato per categoria di prodotto;
  \item ordinare il catalogo per nome, prezzo o tipo;
  \item visualizzare la pagina di dettaglio di un prodotto;
  \item cercare prodotti per nome;
  \item aggiungere e rimuovere prodotti dal carrello;
  \item registrarsi, operazione richiesta solo per concludere un acquisto.
\end{itemize}

\subsection{Utente registrato}

Oltre a quanto sopra:

\begin{itemize}
  \item concludere un acquisto;
  \item consultare lo storico dei propri ordini;
  \item richiedere e scaricare la fattura di un ordine;
  \item modificare i propri dati: nome, cognome, indirizzo predefinito e metodo
        di pagamento.
\end{itemize}

\subsection{Amministratore}

\begin{itemize}
  \item inserire, modificare ed eliminare i prodotti del catalogo tramite
        un'interfaccia dedicata, senza richiedere competenze di programmazione;
  \item consultare tutti gli ordini, filtrandoli per cliente e per intervallo di
        date.
\end{itemize}

\section{Diagramma navigazionale}

\begin{figure}[H]
  \centering
  \resizebox{\linewidth}{!}{\diagrammaNavigazionale}
  \caption{Diagramma navigazionale del sito.}\label{fig:navigazionale}
\end{figure}

\section{Mappa dei contenuti}

\begin{figure}[H]
  \centering
  \resizebox{\linewidth}{!}{\mappaContenuti}
  \caption{Mappa dei contenuti: i collegamenti percorribili fra le pagine.
           La codifica cromatica è quella della figura~\ref{fig:navigazionale}.}
  \label{fig:contenuti}
\end{figure}

\section{Base di dati}

Lo schema concettuale è organizzato attorno a quattro concetti: l'utente, che si
specializza in cliente e amministratore; l'ordine, eseguito da un cliente; il
prodotto, a sua volta specializzato nelle tre categorie in vendita; e l'acquisto,
che associa un ordine ai prodotti che lo compongono registrandone quantità,
prezzo e aliquota IVA al momento dell'acquisto.

\begin{figure}[H]
  \centering
  \schemaER
  \caption{Schema concettuale della base di dati.}\label{fig:er}
\end{figure}

\section{Layout}

\begin{figure}[H]
  \centering
  \resizebox{0.78\linewidth}{!}{\wireframeHome}
  \caption{Wireframe della pagina principale.}\label{fig:layout}
\end{figure}

\section{Scelta dei colori}

La palette è costruita attorno a un blu primario, impiegato per gli elementi
interattivi e per gli accenti, affiancato da una scala di neutri per gli sfondi e
per il testo.

\begin{figure}[H]
  \centering
  \paletteColori
  \caption{Palette cromatica del sito.}\label{fig:colori}
\end{figure}

\begin{table}[H]
  \centering
  \small
  \begin{tabular}{llp{7cm}}
    \toprule
    \textbf{Colore} & \textbf{Codice} & \textbf{Impiego} \\
    \midrule
    Deep Sky Blue & \texttt{\#007BFF} & Colore primario: titoli, collegamenti ed elementi in evidenza \\
    Blue Eyes     & \texttt{\#1663C7} & Variante scura del primario, per gli stati attivi \\
    Snow Drift    & \texttt{\#F8F9FA} & Sfondo delle pagine \\
    White         & \texttt{\#FFFFFF} & Sfondo delle schede prodotto e dei riquadri \\
    Dune          & \texttt{\#333333} & Testo, barra di navigazione e footer \\
    \bottomrule
  \end{tabular}
  \caption{Impiego dei colori nell'interfaccia.}\label{tab:colori}
\end{table}

\clearpage
\section{Il sistema realizzato}

Questa sezione descrive sinteticamente l'applicazione effettivamente sviluppata
e segnala i punti in cui si discosta dalla proposta.

\subsection{Architettura}

L'applicazione è scritta in Java secondo il modello \emph{Model--View--Controller}:
le servlet fanno da controller, le pagine JSP da viste e la persistenza è
affidata al pattern DAO su MySQL tramite JDBC. Il pool di connessioni è dichiarato
come risorsa JNDI e gestito dal contenitore, non dall'applicazione.

\begin{table}[H]
  \centering
  \small
  \begin{tabular}{llp{6.4cm}}
    \toprule
    \textbf{Componente} & \textbf{Versione} & \textbf{Ruolo} \\
    \midrule
    Java SE             & 11 & Livello linguistico di compilazione \\
    Servlet API         & 4.0 (\texttt{javax.servlet}) & Servlet, filtri e listener \\
    Apache Tomcat       & 9.0 & Contenitore web e pool di connessioni \\
    MySQL               & 8.0 & Base di dati \\
    Apache PDFBox       & 2.0.29 & Generazione delle fatture \\
    Google Gson         & 2.11.0 & Serializzazione JSON per gli endpoint asincroni \\
    Apache Maven        & 3.9 & Compilazione e produzione dell'archivio WAR \\
    \bottomrule
  \end{tabular}
  \caption{Tecnologie impiegate.}\label{tab:tecnologie}
\end{table}

Tomcat~9 è un vincolo e non una preferenza: il codice usa il namespace
\texttt{javax.servlet}, che dalla versione~10 è stato rinominato in
\texttt{jakarta.servlet}.

Il codice è organizzato in quattro package: \texttt{control} (21 servlet),
\texttt{model} (bean e DAO), \texttt{filter} (controllo degli accessi e codifica
dei caratteri) e \texttt{config} (configurazione risolta all'avvio). Nessuna
pagina JSP accede direttamente alla base di dati: riceve dalla servlet i dati già
pronti come attributi della richiesta o della sessione.

\subsection{Servlet, viste e filtri}

Il controllo degli accessi non è ripetuto in ogni servlet ma concentrato in tre
filtri, a cui se ne aggiunge un quarto per la codifica dei caratteri.

\begin{table}[H]
  \centering
  \small
  \begin{tabular}{llp{6.6cm}}
    \toprule
    \textbf{Filtro} & \textbf{Pattern} & \textbf{Comportamento} \\
    \midrule
    \texttt{EncodingFilter} & \texttt{/*} & Impone UTF-8 sul corpo delle richieste \\
    \texttt{UserFilter} & \texttt{/user/*} & Richiede un utente autenticato \\
    \texttt{AdminFilter} & \texttt{/admin/*} & Richiede il ruolo di amministratore \\
    \texttt{JspFilter} & \texttt{/pages/views/*} & Nega ogni accesso diretto alle viste \\
    \bottomrule
  \end{tabular}
  \caption{Filtri e pattern di URL.}\label{tab:filtri}
\end{table}

\texttt{JspFilter} merita una nota: le pagine sotto \texttt{pages/views} sono
viste e non risorse pubbliche, e senza i dati preparati da una servlet
produrrebbero errori. Il filtro nega quindi ogni richiesta diretta, ma le pagine
restano raggiungibili tramite l'inoltro interno, che il dispatcher predefinito
non intercetta.

\begin{landscape}
\begin{figure}[H]
  \centering
  \resizebox{!}{13.4cm}{\mappaServlet}
  \caption{Mappa delle servlet, delle viste e delle aree protette dai filtri.
           In azzurro le servlet, in giallo le viste JSP, in rosso le risposte
           generate dinamicamente, in verde il filtro sulle viste.}
  \label{fig:servlet}
\end{figure}
\end{landscape}

\subsection{Dallo schema concettuale allo schema realizzato}

Lo schema effettivo è più semplice di quello di figura~\ref{fig:er}, per due
scelte prese in fase di sviluppo:

\begin{itemize}
  \item \textbf{Cliente} e \textbf{Admin} non sono tabelle distinte: la
        specializzazione è resa da un campo \texttt{tipo} nella tabella
        \texttt{Utente}, che assume i valori \texttt{user} e \texttt{admin};
  \item le specializzazioni di \textbf{Prodotto} non hanno attributi propri, per
        cui sono anch'esse ridotte a un campo \texttt{tipo}
        (\texttt{keyboard}, \texttt{switch}, \texttt{keycap}).
\end{itemize}

Restano quindi quattro tabelle collegate in cascata: un \texttt{Utente} ha più
\texttt{Ordine}, ogni ordine ha più righe di \texttt{Acquisto}, e ogni riga
riferisce un \texttt{Prodotto}.

Tre scelte progettuali sono degne di nota:

\begin{description}
  \item[Dati fotografati all'acquisto.] \texttt{Acquisto} non si limita a
    collegare ordine e prodotto: ne copia prezzo, aliquota IVA e nome del file
    immagine. Una variazione di listino non altera così gli ordini già conclusi,
    e una fattura ristampata a distanza di mesi riporta gli importi corretti.
  \item[Indirizzo per ordine.] L'indirizzo di consegna è memorizzato su ogni
    \texttt{Ordine} e non letto dal profilo, perché lo stesso cliente può
    spedire a destinazioni diverse. Quello del profilo resta il valore proposto
    per impostazione predefinita.
  \item[Cancellazione logica.] L'eliminazione di un prodotto imposta
    \texttt{tipo = 'Eliminato'} anziché rimuovere la riga: una cancellazione
    fisica violerebbe il vincolo di chiave esterna con \texttt{Acquisto} e
    cancellerebbe lo storico degli ordini.
\end{description}

\subsection{Compilazione ed esecuzione}

Il progetto nasceva come \emph{Dynamic Web Project} di Eclipse: la configurazione
del server, del percorso delle immagini e delle credenziali viveva nel workspace
dell'IDE e non nel repository, rendendo il progetto non riproducibile altrove.

La compilazione è ora affidata a Maven, che produce l'archivio WAR e risolve le
librerie in luogo dei file JAR precedentemente versionati. L'esecuzione è
affidata a Docker Compose, che avvia due contenitori: MySQL, che alla prima
creazione del volume esegue lo schema e i dati di esempio, e Tomcat~9 con
l'applicazione, costruita da un \texttt{Dockerfile} a due fasi. Un solo comando
avvia l'intero sistema, senza installare Maven, MySQL o Tomcat.

Le credenziali non sono più scritte nel repository: \texttt{context.xml} contiene
segnaposto che Tomcat risolve all'avvio leggendo le variabili d'ambiente, oppure
le \emph{system property} quando l'applicazione viene avviata da IDE.

\begin{lstlisting}[language=XML,caption={Dichiarazione del data source.},label=lst:context]
<Resource name="jdbc/keyItaly"
          auth="Container"
          type="javax.sql.DataSource"
          driverClassName="com.mysql.cj.jdbc.Driver"
          url="${DB_URL}"
          username="${DB_USERNAME}"
          password="${DB_PASSWORD}"
          connectionProperties="useUnicode=true;characterEncoding=UTF-8"/>
\end{lstlisting}

Analogamente, la directory delle immagini del catalogo --- esterna
all'applicazione perché scrivibile dall'amministratore --- viene individuata
all'avvio da un listener che la cerca fra variabile d'ambiente, system property e
parametro di contesto, creandola se assente. Se non è utilizzabile
l'applicazione non parte: l'errore diventa così immediatamente visibile invece di
manifestarsi come una serie di immagini mancanti.

\subsection{Limiti noti}

Il progetto ha finalità didattiche e alcune scelte, accettabili in quel contesto,
non sarebbero adeguate a un sistema in esercizio.

\begin{itemize}
  \item \textbf{Password in chiaro}, senza funzione di hash.
  \item \textbf{Assenza di protezione CSRF}; alcune operazioni di modifica
        accettano anche richieste GET.
  \item \textbf{Output non filtrato}: le JSP scrivono i valori senza codifica
        delle entità HTML, con possibile cross-site scripting persistente.
  \item \textbf{Giacenza non verificata}: l'acquisto non controlla né decrementa
        la disponibilità.
  \item \textbf{Nessun test automatico}: la verifica è stata condotta
        manualmente.
\end{itemize}

Le evoluzioni naturali sono l'introduzione dell'hashing delle password e dei
token CSRF, la sostituzione degli scriptlet con JSTL ed \emph{expression
language} --- che darebbe la codifica automatica dell'output --- e la gestione
transazionale della conclusione dell'ordine.

\end{document}
